% Pursuing Stacks (cotes 134-2 to 134-8): archive cote 134-2, site page 28 (PDF page 29).
% letter to Quillen. Draft compared: ps1.tex:660-736.
% AI-transcribed by Claude Opus 5.5 (claude-opus-5-5, adaptive thinking, effort high), Message Batches API, batch ps1 (msgbatch_01Rv5hCaNTtUt4JHx1yqoHPL), 2026-11-12; not yet checked by eye.
\scanpage{134-2/28}
\margin{8}
\mspage{\deleted{8} \hand{14}}
21.2.

\over{\unsure{4}}{8}. Returning to the\unsure{amalgamated|amalgamated} sum $X = \hand{\bigcup_i}X_i$ , the cellular decompositions
of the components $X_i \simeq D_{n(i)}$ define a cellular decomposition of X, whose
set of cells with incidence relation forms a finite ordered set K, fini-%
te union of a family of subsets $(K_i)_{i\in I}$ , with given isomorphisms\xout{m}
\[ f_i : K_i \simeq \underline{J}_{n(i)} \quad (i\in I) , \]
where for every index $n\in \underline{N}$, $\underline{J}_n$ denotes the ordered set of the 2n+1
cells $S_j^{\hand{\varepsilon}}$ $(0\leq j\leq n-1,\ \hand{\varepsilon\in\{}\pm 1\hand{\}})$, $D_n$ of the \deleted{pertinent} \inserted{relevant} cellular decomposition
of $D_n$ . We may without loss of generality assume there is no\deleted{t} inclusion
relation between the $K_i$ , moreover the standar\unsure{\over{t}{d}}ness condition described
above readily translates into a condition on this structure
$K$, $(K_i)_{i\in I}$ , $(f_i)_{i\in I}$ , and implies that for $i,i'\in I$, $K_i\hand{\cap}K_{i'}$ is a``clo-%
sed'' subset in the two ordered sets $K_i,K_{i'}$ (namely contains with any
element x the elements smaller than x), and moreover isomorphic (for the
induced order) to some $\underline{J}_n$ . Thus the category $B_o$ \xout{is} can be viewed as the
category of such ``standar\unsure{\over{t}{d}} ordered sets'' (with the extra structure on the-%
se just said), and the category $C_o$ can be defined most simply as the
dual category $B_o^{\,o}$ . (NB the definition of morphisms in $B_o$ is clear I
guess ...) I believe the category $B_o$ is stable under amalgamated sums
$X\amalg_Z Y$ , provided \inserted{however} we insist that the empty structure K is \uline{not} allowed -
otherwise we have to \xout{\unsure{insist}} restrict to amalgamated sums with $Z \neq \emptyset$ .
It seems finally more convenient to exclude the empty structure in $B_o$ ,
\margin{\unsure{(13!)}}
i.e. to exclude final element from $\underline{C}_o$ , for the benefit of being able to
\margin{* This is \unsure{obviously} false, see P.S. p.12}
state that $\underline{C}_o$ (and \uline{all categories $\underline{C}_n$}\hand{*} \deleted{\illegible{1}} \xout{\illegible{1}} \margin{(13!)} are stable under binary
amalgamated sums, and that the functors $\underline{C}_n\to C$ (and ultimately $\underline{C}_\infty\to C$)
we are interested in are just those commuting to arbitrary binary amalga-%
mated sums (without \xout{\illegible{1}} awkward reference to the objects $\underline{F}_i$ and the
iterated source and target maps between them ...)\hand{(13!)}.

\over{\illegible{1}}{9}. The category $\underline{C}_o$ being fairly well understood, \inserted{it} remains to complete
the construction by the inductive step, \inserted{\unsure{for)}} passing from $\underline{C}_n$ to $\underline{C}_{n+1}$ . The
main properties I have in mind \deleted{therefore} \inserted{for this}, for the sequence of categories
$\underline{C}_n$ and their limit $\underline{C}_\infty$ , are the following two.

A) For any $K\in Ob(\underline{C}_\infty)$ \ ($= Ob(\underline{C}_o)$), and any two arrows in $\underline{C}_\infty$
\[ f,g : K \rightrightarrows F_i , \]
with $i\in\underline{N}$, \hand{(}and such that either i=0, or the equalities
\[ \hand{(1)}\qquad s_1^i f = s_1^i g \ , \quad t_1^i f = t_1^i g \]
hold (case $i \over{>}{\geq} 1$), there exists $h: K\to F_{i+1}$ such that
\[ \hand{(2)}\qquad s_1^{i+1} h = f \ , \quad t_1^{i+1} h = g \ . \]
B) For any $n\in\underline{N}$, the category $\underline{C}_{n+1}$ is deduced from $\underline{C}_n$ by keeping
the same objects, and just adding new arrows \xout{\illegible{1}} h as in A), with
f, g arrows in $\underline{C}_n$ .

\margin{28}
